\section{A warning in a case note is not the effect of a policy}
\label{closure30:text-signals}

A bank reading thousands of case notes faces two different questions. Which
borrowers appear to be in difficulty? Would a change in the terms offered to
them reduce that difficulty? A language model can help answer the first and
organise evidence for the second. It can also make the distinction harder to
see. If borrowers receiving relief are already in worse condition, a model
may correctly associate the words describing relief with subsequent default.
The association can be large even when relief helps.

Three studies in pharmacovigilance make this distinction concrete. Their
subject is medical reporting, not banking, but the proposal faces the same
question about what a text-derived comparison means. InferBERT combines an
ALBERT classifier with comparisons between reports containing and not
containing a feature. DeepCausality extends the approach to named entities
in free text. A later assessment compares general, biomedical and
task-specific BERT models. These are useful sources for text processing and
candidate discovery. Their intervention interpretation requires a separate
examination of the actual calculation.
\citep{closure30inferbert,closure30deepcausality,closure30pharma}

\subsection{What the primary methods compare}

Let $W$ be a report, $Y$ its outcome label and $E(W)$ indicate the presence of
a candidate term. Write $m(W)$ for the classifier's predicted probability of
$Y=1$. The root comparison described in InferBERT's equations (2)--(4) and
DeepCausality's named-entity algorithm is, in population notation,
\begin{equation}
 G_E=\E\{m(W)\mid E(W)=1,S=1\}
       -\E\{m(W)\mid E(W)=0,S=1\},
 \label{closure30:entity-gap}
\end{equation}
where $S=1$ denotes the eligible, nonempty records retained for that feature.
The corresponding source code reads already computed classifier
probabilities, divides record indices into two groups according to feature
presence, and compares group means. Its statistic has the form
\begin{equation}
 Z_E=\frac{\overline m_1-\overline m_0}
 {\sqrt{s_1^2/n_1+s_0^2/n_0}},
 \label{closure30:entity-z}
\end{equation}
with sample variances $s_e^2$ and group sizes $n_e$. The inspected routines
do not replace a structural equation governing a person's treatment or
recompute that person's outcome under an alternative treatment. The papers
call the procedure do-calculus; equations~\eqref{closure30:entity-gap} and
\eqref{closure30:entity-z} describe the implemented group comparison.
\citep[pp.~5--6, equations (2)--(6)]{closure30inferbert}
\citep[pp.~4--6, Algorithm 1]{closure30deepcausality}
\citep{closure30inferbertcode,closure30deepcausalitycode}

The distinction survives a very good classifier. If
$m(W)=\E(Y\mid W)$ exactly and both $E$ and eligibility $S$ are functions
of the recorded information, iterated expectation gives
\begin{equation}
 G_E=\E(Y\mid E=1,S=1)-\E(Y\mid E=0,S=1).
 \label{closure30:calibration-group}
\end{equation}
This is a difference between selected populations. It need not equal the
effect of changing treatment in a common population. A small standard error
in~\eqref{closure30:entity-z} can establish that the populations differ while
leaving that identification question unanswered. General calibration by
score alone is weaker than the conditional prediction equality used here;
the argument deliberately grants the classifier the stronger property.

There are also limits to what can be reproduced from the inspected revisions.
The InferBERT scoring routine filters one-sided normal-reference probabilities
at 0.05. In the public DeepCausality notebook that filter is commented out,
although the paper discusses adjusted probabilities. The notebook still
calculates and exports probabilities and ranks terms. That difference does
not prove that the published numerical results are wrong: further analysis
may have occurred elsewhere. It does mean that this notebook alone does not
document the paper's complete multiplicity procedure. The two revisions are
fixed at their May 2021 and November 2022 commits, respectively.
\citep{closure30inferbertcode,closure30deepcausalitycode}

\subsection{What the reported results establish}

InferBERT constructs two FAERS-based tasks, using 36,661 analgesics reports
and 27,245 tramadol reports after processing. Its upstream source is the
commercially curated PharmaPendium database. The study uses a stratified
64:16:20 training, development and test split. Thus public availability of
FAERS in general is not evidence that an independently prepared public
extract reproduces this particular curated cohort.
\citep[pp.~3--4, Table 1]{closure30inferbert}

DeepCausality constructs 14,361 sentences from LiverTox, carrying the
drug-level likelihood classification into sentence labels. It reports a
stratified 90:10 sentence split and selects an optimised model using the
reported test-set classification results. That evaluation should not be
described as an untouched held-out assessment on new drugs or new patients.
Nor does the split alone establish that a particular amount of leakage
occurred. The paper reports that 18 of 20 enriched terms agree with clinical
knowledge; it also acknowledges that the method cannot distinguish
confounders from colliders. These statements support a role in finding
relevant concepts, while delimiting a role in selecting adjustment variables.
\citep[pp.~5--8, Table 1; pp.~9--11, Table 2 and Discussion]
{closure30deepcausality}

The 2023 assessment offers a useful warning against choosing a model solely
by classification accuracy. In its LiverTox task, several models report an
accuracy of 0.91, while the expert-reviewed relevance of enriched terms
differs. The reported positive predictive value for DILI-BERT is 0.95,
compared with 0.20 for RoBERTa and 0.25 for SciBERT. The paper's quantity
labelled predictive positive rate is $TP/(TP+FP)$; Figure 5 concerns the
relevance of enriched terms. None of these numbers is an estimate of the
effect of prescribing a drug, granting a loan concession, or changing a
collection policy. They are reported source results, not independently
reproduced performance in this proposal.
\citep[pp.~1911--1916, Table 4 and Figure 5]{closure30pharma}

The proposal should therefore retain both lessons. Domain-specific training
can change which concepts a model ranks as relevant, even when its headline
classification score changes little. Expert relevance, repeatability and
prediction remain different from a valid intervention contrast. Rejecting
the latter interpretation does not erase the former contribution.

\subsection{A lending example with a perfectly calibrated text model}

Consider a deliberately constructed hardship programme. Let $U$ denote an
unrecorded high-risk condition, $A$ receipt of relief, $R$ a noisy baseline
description of that condition, and $Y$ subsequent default. All are binary.
Let $U$ have probability one half; conditional on $U$, independently generate
\begin{equation}
 \Pr(A=U\mid U)=\frac45,\qquad
 \Pr(R=U\mid U)=\frac34.
 \label{closure30:text-design}
\end{equation}
The outcome mechanism is
\begin{equation}
 \Pr(Y=1\mid U,A)=\frac1{10}+\frac35U-\tau A,
 \qquad \tau\in\left\{0,\frac1{20}\right\}.
 \label{closure30:text-outcome}
\end{equation}
An independent uniform disturbance generates this Bernoulli outcome.
Relief is more common among high-risk borrowers. It either has no effect,
or reduces default probability by five percentage points. These probabilities
illustrate a mechanism; they are not estimates for a bank.

Suppose the document truthfully records $W=(A,R)$, and the language model
returns the exact conditional probability $m(a,r)=\Pr(Y=1\mid A=a,R=r)$.
Figure~\ref{closure30:text-dag} distinguishes the economic variables from
their description. Rewriting a note is an operation on its representation;
granting relief changes the economic variable $A$.

\begin{figure}[htbp]
\centering
\includegraphics[width=.89\textwidth]{twelve-hour-closure/2026-10-29/text-figures/decision.pdf}
\caption{The constructed hardship example. Baseline risk influences relief,
the noisy description and default. The $A\to Y$ link has coefficient
$-\tau$ in the conditional default probability and is absent when $\tau=0$.
The box records the text available to the predictor; it is not an additional
cause of the customer's default in this example.}
\label{closure30:text-dag}
\end{figure}

Bayes' rule gives, in the order $(a,r)=(0,0),(0,1),(1,0),(1,1)$,
\begin{equation}
 \Pr(U=1\mid A=a,R=r)
   =\left(\frac1{13},\frac37,\frac47,\frac{12}{13}\right).
 \label{closure30:text-posterior}
\end{equation}
For $\tau=0$, the four predictions are therefore
$(19/130,5/14,31/70,17/26)$. For $\tau=1/20$, subtract $1/20$ from the
two predictions with $a=1$. Every prediction is conditionally calibrated.
Nevertheless the entity-group contrast is
\begin{equation}
 G_A=\frac35\left(\frac45-\frac15\right)-\tau
     =\frac9{25}-\tau.
 \label{closure30:text-observed}
\end{equation}
It is positive in both cases. A source-style test can increasingly detect
this group difference as independent observations accumulate. The question
it answers is still a group comparison.

Now consider the superficially stronger operation of rewriting the relief
term in each document while leaving $R$ fixed. This is an additional
illustrative comparison, not a claim that the inspected source routine
performs such rewriting. Averaging the predictor changes gives
\begin{equation}
 G_{\mathrm{rewrite}}
 =\E_R\{m(1,R)-m(0,R)\}=\frac{27}{91}-\tau.
 \label{closure30:text-rewrite}
\end{equation}
All four document combinations have positive probability, so the example
does not depend on predictions for impossible text. The predictor changes
its inferred distribution of $U$ when the relief term changes. An actual
intervention leaves the original distribution of $U$ in place and replaces
the treatment assignment. Hence
\begin{equation}
 \E\{Y\mid\operatorname{do}(A=1)\}
 -\E\{Y\mid\operatorname{do}(A=0)\}=-\tau.
 \label{closure30:text-do}
\end{equation}

\begin{table}[htbp]
\centering
\begin{tabular}{lrr}
\toprule
Comparison & No relief effect & Relief lowers default\\
\midrule
Difference between text-score groups & $36.0$ & $31.0$\\
Change in score after text rewriting & $29.67$ & $24.67$\\
Actual effect of relief & $0.0$ & $-5.0$\\
Adjustment using truthful baseline $U$ & $0.0$ & $-5.0$\\
\bottomrule
\end{tabular}
\caption{Exact population differences, expressed in percentage points, in
the two constructed economies. The predictions are conditionally calibrated
in both. The values follow from exact enumeration of all 16 joint cells and
the independent formulas in the text; they are not estimates from bank data.}
\label{closure30:text-results}
\end{table}

The example identifies a practical danger. Rejecting relief because its
mention predicts default would reject a beneficial policy in the second
economy. Better text prediction does not resolve the comparison, because
the prediction itself is already correct. The missing information concerns
why the policy was assigned.

\subsection{A constructive route: text as measured baseline information}

The appropriate conclusion is conditional adoption, rather than abandoning
text. In the constructed economy, observing truthful baseline $U$ would
justify comparing relief within its two levels and then averaging over the
common baseline population. The resulting effect is exactly $-\tau$.
Randomising relief would provide a different route by breaking its relation
with baseline risk. Neither route is produced merely by relabelling a
classifier contrast as an intervention.

There is a constructive literature on preserving adjustment information
while reducing text to a manageable representation. \citet{closure30veitch}
begin by assuming that the text contains enough information for adjustment.
Their Theorem 3.1 gives conditions under which an embedding preserves the
relevant conditional outcome or propensity information; Theorem 3.2 adds
estimation conditions. The initial sufficiency assumption is essential.
An embedding cannot acquire an omitted baseline cause simply by predicting
the observed outcome accurately. For the hardship example, a baseline
narrative would need a defensible relationship to $U$; the noisy signal $R$
alone does not become sufficient by being written in fluent language.

The target population must also remain explicit. Veitch and colleagues
develop an effect on the treated, with equation (2.1) averaging over treated
units. The present constructed comparison is an average effect in the full
population. A bank must select its target before choosing the weights in an
estimator. The source's theorem display omits the treated conditioning in
its final expectation while the earlier definition and estimators retain it;
the proposal uses the explicit target definition rather than silently
changing the population.
\citep[Section 2, equations (2.1)--(2.3); Section 3, Theorems 3.1--3.2]
{closure30veitch}

For a practical evaluation, the bank would first fix the decision time and
separate baseline information from later descriptions of treatment and
outcome. It would then assess extraction accuracy, agreement with audited
baseline facts, the assumptions supporting adjustment, overlap between
treatment groups, and effect uncertainty. Borrower- or document-grouped
splits and a final unused evaluation set address predictive evaluation;
they do not by themselves supply missing confounding information. Public
documents and generated assignment laws can test these distinctions before
an expensive institutional experiment is justified.

\subsection{What to adopt, integrate and decline}

The reusable idea is a traceable route from documents to candidate
variables and ranked associations. Adopt that conceptual separation.
Integrate the extracted term, supporting passage, information date,
eligibility rule and score as evidence for a reviewer. Preserve whether a
term denotes a treatment, a baseline cause, a symptom, a consequence or an
uncertain relation. A relevant term should not automatically become an
adjustment variable or a directed causal arrow.

Direct reuse of the inspected implementations requires further work. Both
repositories target TensorFlow 1.15; DeepCausality specifies Python 3.6 and
external BioBERT weights and BERN results. The repositories refer to external
data downloads, and InferBERT's cohort derives from commercial curation.
The inspected InferBERT tree has no repository-wide licence file. In
DeepCausality, an Apache licence accompanies the embedded BioBERT subtree;
the project-specific scoring notebook has no explicit licence notice.
These observations leave the permission for copying the project-specific
code unresolved; a dependency's notice does not answer that question.
No legal clearance or native reproduction is claimed here.
\citep{closure30inferbertcode,closure30deepcausalitycode}

For the proposal, favour an independently specified, tested implementation
of the disclosed mathematical comparison if that diagnostic is useful and
its permitted scope is settled. Keep it labelled as an association or
hypothesis-ranking result. Decline its direct promotion to a general
intervention estimator. A separate effect calculation should require a
defined policy, time ordering, identification assumptions and appropriate
uncertainty. This preserves the value of language assistance while preventing
a familiar business error: mistaking the customers selected for a policy
for the customers' response to it.
